
\vspace{0.2cm}

\subsection{Architectural Pattern Selection}

The automated data extraction pipeline, defined in the BD5 task in Section \ref{sec:taskdescription}, can be implemented using two distinct architectural patterns. The first is a synchronous monolithic architecture, in which the entire process, from the initial upload to the AI processing of the expense, is handled through a single API call. The second is an asynchronous microservices-based architecture, where each service is responsible for a specific function in the workflow.

An expense upload triggers the POST /api/expenses call. A simple POST request is quick and lightweight, whereas the AI-powered extraction introduces considerable processing delay. Consequently, the two architectural patterns exhibit opposite behaviours.

In the case of synchronous monolithic architecture, the POST call cannot complete until AI extraction is finished, so the server response time is considerably prolonged. This is particularly harmful in bulk-upload scenarios, where expenses are processed sequentially, degrading user experience and reducing server availability. At the same time, reliability is low: any error in the end-to-end workflow causes the entire request to fail, requiring users to retry the upload.

In an asynchronous microservices-based architecture, the POST call is fully decoupled from the AI extraction process. Single uploads complete instantly, and bulk uploads are handled concurrently, preventing server overload. After a successful expense upload, the expense is routed to a queue and processed by the AI when resources are available. In addition, a more robust error-handling strategy is applied: if the AI model fails to extract the data, the expense is preserved and sent back to the queue for a later retry.

Due to its negative impact on user experience and scalability, the synchronous monolithic approach is discarded, biasing the design toward the alternative solution. The \textbf{asynchronous architecture with microservices} delivers immediate feedback to the user, prevents API bottlenecks, and results in a more resilient and scalable system. By making this choice, the non-functional requirements NFR1, NFR3, and NFR4 (Section \ref{sec:non-functional-requirements}) are also fulfilled.
